% !TeX root = ../初等数学笔记.tex
% 由主文档通过 \input 载入；请编译主文档。
\setcounter{chapter}{10} % 保留原书章号。
\chapter{计数原理与二项式定理}
\label{chap:em-counting}
计数的首要工作是准确规定对象：顺序是否重要、同一元素能否重复、哪些表示被视为同一个结果。公式只有在这些约定明确之后才有意义。

\section{加法原理、乘法原理与对应}
若有限集合 $A_1,\ldots,A_m$ 两两不交，则其并集大小为 $\sum_i|A_i|$，这称为加法原理。若一个选择过程分为 $m$ 步，且无论前面如何合法选择，第 $i$ 步都有恰好 $n_i$ 种选择，则总数为 $\prod_i n_i$，这称为乘法原理；若分支数不同，应分别计算再相加。

两个集合之间若能建立一一对应的\emph{双射}（bijection），则它们元素个数相同。例如从 $n$ 个对象中选出一个子集，等价于给每个对象指定选或不选，因此共有 $2^n$ 个子集。这里计数的是子集，而不是选择对象的先后过程。

\section{排列、组合与阶乘}
规定 $0!=1$，正整数 $n$ 的阶乘为 $n!=1\cdot2\cdots n$。从 $n$ 个不同对象中不重复地依次选 $k$ 个，$0\le k\le n$，有 $n(n-1)\cdots(n-k+1)=n!/(n-k)!$ 种。若不考虑顺序，每个 $k$ 元子集对应 $k!$ 个有序结果，所以\emph{组合数}（binomial coefficient）为
\[
 \binom nk=\frac{n!}{k!(n-k)!}.
\]
$\binom n0=1$ 对应唯一的空选择。

若 $n$ 个位置中有 $m_1,\ldots,m_r$ 个分别相同的对象，总数 $\sum_i m_i=n$，则不同排列数为 $n!/(m_1!\cdots m_r!)$；分母消去的是每个结果中相同的重复标记次数。$n\ge1$ 个不同对象围成圆环，若只把旋转视为相同而反射视为不同，固定一个对象为起点后有 $(n-1)!$ 种。改变等价关系会改变答案，不能再无条件除以 $2$。

\section{重复选择、分配与双阶乘}
非负整数解 $x_1+\cdots+x_k=n$（$n\ge0,k\ge1$）可用 $n$ 个星号与 $k-1$ 个隔板编码，隔板前后星号数量给出各 $x_i$，因而共有 $\binom{n+k-1}{k-1}$ 个。若要求每个 $x_i\ge1$，先令 $y_i=x_i-1$；当 $n\ge k$ 时共有 $\binom{n-1}{k-1}$ 个，否则无解。上界限制还需进一步筛选或容斥。

双阶乘每次递减 $2$ 相乘，约定 $0!!=1$。对 $n\ge1$，
\[
 (2n)!!=2^n n!,\qquad
 (2n-1)!!=\frac{(2n)!}{2^n n!}.
\]
将 $2n$ 个不同对象两两配对，固定第一个对象后它有 $2n-1$ 个搭档，再对剩余对象递归配对，故总数为 $(2n-1)!!$。也可从全排列中除以每对内部交换的 $2^n$ 与各对重新排序的 $n!$，得到同一公式。

\section{二项式定理与组合恒等式}
\begin{theorem}[二项式定理]
对非负整数 $n$，
\begin{equation}\label{eq:em-binomial}
 (x+y)^n=\sum_{k=0}^n\binom nk x^{n-k}y^k.
\end{equation}
\end{theorem}
\begin{proof}
当 $n\ge1$，展开 $n$ 个因子的乘积，每一项从每个因子中选 $x$ 或 $y$。恰选 $k$ 次 $y$ 的选择数为 $\binom nk$，于是得到对应系数。$n=0$ 时等式两端均为常数多项式 $1$，按多项式零次幂约定理解。
\end{proof}
把 $k$ 元子集按是否含指定元素分类，得到 Pascal 递推 $\binom nk=\binom{n-1}k+\binom{n-1}{k-1}$，边界之外的组合数约定为 $0$。取补集给出 $\binom nk=\binom n{n-k}$，式 \eqref{eq:em-binomial} 取 $x=y=1$ 得 $\sum_k\binom nk=2^n$；取 $x=1,y=-1$ 可比较偶数大小与奇数大小子集的数目。

\section{容斥与不重不漏的计数}
对有限集合 $A,B$，
\[
 |A\cup B|=|A|+|B|-|A\cap B|,
\]
因为交集中的元素在前两项被计了两次。三个集合相应满足
\[
 |A\cup B\cup C|=|A|+|B|+|C|
 -|A\cap B|-|A\cap C|-|B\cap C|+|A\cap B\cap C|.
\]
一般的\emph{容斥原理}（inclusion--exclusion principle）按交集大小交替加减。一个属于恰好 $r\ge1$ 个集合的元素得到的总系数为 $\sum_{j=1}^r(-1)^{j+1}\binom rj=1$，由二项式定理可知，因此每个并集元素恰计一次。

例如 $1$ 至 $100$ 中被 $2$ 或 $3$ 整除的整数有 $50+33-16=67$ 个，被两者整除等价于被 $6$ 整除。若要计算两者都不整除的数量，用全集大小减去并集即可，答案为 $33$。

\section{例题、习题与学习检查}
\begin{example}
把 $8$ 个相同物品分给 $3$ 个不同的人，允许有人分不到，且第一人至多分到 $3$ 个，共有多少种分法？
\end{example}
\begin{solve}
不限制上界时有 $\binom{10}2=45$ 种。违反限制意味着第一人至少有 $4$ 个，先给他取走 $4$ 个，余下 $4$ 个仍分给三人，共 $\binom62=15$ 种。故合法分法为 $45-15=30$ 种。相同物品没有内部排列，不能乘以 $8!$。
\end{solve}
\begin{enumerate}
\item 从 $6$ 人中选 $3$ 人并指定其中一人为负责人。答：$\binom63\cdot3=60$。
\item 将 $6$ 个不同对象两两配对。答：$5!!=15$。
\item 证明 $\sum_{k=0}^n k\binom nk=n2^{n-1}$（$n\ge1$）。提示：计数带有一个指定成员的非空子集，先选该成员再选其余成员。
\end{enumerate}

